\section{Occupation holonomy, finite resonances and their endpoint residues}
\label{sec:v40-holonomy-residues}
The new full-band estimate permits the measurable phase group to be
used as an actual peripheral spectrum. We first strengthen its
classification. The occupation term changes a stable or unstable
holonomy by an integer. Countability of those changes, rather than
continuity of an occupation itinerary, excludes every nonzero roof
coordinate.

\begin{theorem}[No roof coordinate in an occupation resonance]
\label{thm:v40-no-roof-resonance}
Every solution of \eqref{eq:v39-phase-equation} has $b=0$.
Consequently $\mathcal G_R$ is a finite cyclic group of order
$d_R\le D$, where $D<\infty$ is independent of $R$. Its coordinates
$(u,v,s)$ belong to a fixed finite set of rational torus points, with
possible denominators at most $D$. The lattice alternative in
Theorem~\ref{thm:v39-resonance-structure} is thereby excluded; the
finite cyclic group is not thereby asserted trivial.
\end{theorem}
\begin{proof}
Use a regular positive-measure product block and its stable and
unstable conditional pair measures, as in the measurable holonomy
proof of Theorem~\ref{thm:v35-rotation-free-rigidity}. Restrict the
block if necessary so that the two marginals of each pair measure
are bounded by a constant times $\nu$, and so that leaf contraction
has a common constant. This uses only the local product and absolute
continuity input already used there. For a stable pair $(x,y)$,
$d(T_R^j x,T_R^j y)\le C\Lambda^{-j}$.
The boundary of $Y_R^*$ is a fixed finite collection of rectangle
sides, with uniformly bounded tubular-neighborhood measure. Thus
\[
 \pi^s\{\eta_R(T_R^j x)\ne\eta_R(T_R^j y)\}
       \le C\nu\{d(\cdot,\partial Y_R^*)\le C\Lambda^{-j}\}
       \le C'\Lambda^{-j}.
\]
Invariance of $\nu$ preserves the marginal domination at time $j$.
The first Borel--Cantelli lemma now shows that
\[
 n^s(x,y)=-\sum_{j\ge0}
       (\eta_R(T_R^j y)-\eta_R(T_R^j x))\in\Z
\]
is a finite sum almost everywhere for the stable pair measure.
Backward contraction gives likewise the finite integer
\[
 n^u(x,y)=\sum_{j<0}
       (\eta_R(T_R^j y)-\eta_R(T_R^j x))
\]
for unstable pairs. Independence is not used.

Approximate the measurable circle-valued $q$ by smooth bounded
functions in $L^1$. Marginal domination and leaf contraction show
that the ratio of the two far-future values tends to one in measure
on stable pairs; the far-past ratio does so on unstable pairs.
Fix the approximation first, take time to infinity next, and only
then remove the approximation. The displacement itinerary agrees
along the corresponding regular leaves, while the constant collision
phase cancels between the two endpoints. The action identity and
the finite occupation sums therefore give
\begin{equation}\label{eq:v40-integer-leaf-holonomy}
 \frac{q(y)}{q(x)}
   =\exp\left(ib\int_{[x,y]^{s/u}}\vartheta_R
                     +iv n^{s/u}(x,y)\right)
\end{equation}
almost everywhere for the respective pair measures. The sign in
$n^s$ comes from forward telescoping; $n^u$ uses backward telescoping.

Disintegration and Fubini permit multiplication around almost every
small stable--unstable rectangle $Q$ of this product block. They yield
\begin{equation}\label{eq:v40-countable-area}
 \exp\left(ib\int_Q d\vartheta_R\right)\in
                          \{e^{ivj}:j\in\Z\}.
\end{equation}
The set on the right is countable, including when it is dense in the
circle. The area on the left has no atoms under the product
rectangle parameters. To check this last point, fix three sides and
vary the fourth leaf through a nonatomic transverse conditional
parameter. The resulting rectangles are nested. Their symplectic
area is continuous and strictly monotone, since
$d\vartheta_R=R\cos\varphi\,d\varphi\wedge d\alpha$ has a fixed
nonzero sign on the regular block. Each prescribed area therefore
has at most one value of that transverse parameter. Absolute
continuity of the product coordinates and Fubini give zero measure
to any countable collection of such area levels. The argument also
works on a positive-measure product block with gaps: restrict the
nonatomic transverse parameters to that block.

If $b\ne0$, the inverse image under $A\mapsto e^{ibA}$ of a
countable set is countable. Equation~\eqref{eq:v40-countable-area}
would then hold only on a null set of rectangles, a contradiction.
Hence $b=0$. The uniform cyclic and denominator conclusions now
follow from Theorem~\ref{thm:v39-resonance-structure}. No choice of
periodic representative of $q$ has been made.
\end{proof}

\begin{corollary}[Full-occupation decay away from zero roof frequency]
\label{cor:v40-roof-annulus-gap}
For each $0<\epsilon<B<\infty$ there are $C_{B,\epsilon}<\infty$
and $\rho_{B,\epsilon}<1$ such that
\[
 \|\mathscr T_{R,(u,b,v)}^m\|_B\le C_{B,\epsilon}\rho_{B,\epsilon}^m
 \quad(R\in I,\ u\in\T^2,\ v\in\T,\ \epsilon\le|b|\le B).
\]
\end{corollary}
\begin{proof}
This is a compact joint nonresonant set by the preceding theorem.
Apply Theorem~\ref{thm:v40-full-peripheral}.
\end{proof}

\subsection{The actual source coefficient at every resonance}
Fix $R$ and $\gamma=(u,b,v)\in\mathcal G_R$, and let $s_\gamma$ be
its unique scalar phase. Normalize $|q_\gamma|=1$. Write
$\Pi_\gamma$ for the spectral projection of $\mathscr T_{R,\gamma}$
at $\zeta_\gamma=e^{-is_\gamma}$.

\begin{proposition}[Physical projection and general endpoint residue]
\label{prop:v40-general-residue}
For every bounded physical density $a\nu\in\mathcal B$,
\begin{equation}\label{eq:v40-projection-formula}
 \Pi_\gamma(a\nu)
       =q_\gamma\nu\int\overline{q_\gamma}a\,d\nu.
\end{equation}
For every bounded terminal multiplier $d$ its endpoint coefficient is
\begin{equation}\label{eq:v40-endpoint-residue}
 \ell(M_d\Pi_\gamma(a\nu))
      =\nu(a\overline{q_\gamma})\nu(dq_\gamma).
\end{equation}
In particular the actual section-to-section coefficient is
$A_\gamma=|\nu(\eta_Rq_\gamma)|^2$ at every scalar eigenvalue, not
only at one.

More generally, before using $b=0$, invariance alone gives the exact
identity
\begin{align}\label{eq:v40-residue-source-identity}
 (e^{iv}-1)\nu(\eta_Rq_\gamma)
  ={}&(e^{-is_\gamma}-1)\nu(q_\gamma)\notag\\
     &-\nu\bigl((e^{i(u\cdot\kappa_R+b\tau_R)}-1)
                               e^{iv\eta_R}q_\gamma\bigr).
\end{align}
Thus the pure-occupation cancellation holds when $u=b=s_\gamma=0$
and $v\ne0$. It does not annihilate a general residue with $u\ne0$
or $s_\gamma\ne0$.
\end{proposition}
\begin{proof}
The averages
$M^{-1}\sum_{j<M}\zeta_\gamma^{-j}\mathscr T_{R,\gamma}^j$
converge in the strong space to $\Pi_\gamma$. On physical $L^2$
the same averages converge to the orthogonal projection onto the
unitary eigenspace. That eigenspace is one-dimensional: quotients
of two circle eigenfunctions are invariant and hence constant.
Comparison on smooth global tests and strong-space faithfulness
prove \eqref{eq:v40-projection-formula}. The terminal pairing gives
\eqref{eq:v40-endpoint-residue}. This proof does not require
$M_{q_\gamma}$ or $M_{\overline q_\gamma}$ to be a strong multiplier.

Integrate \eqref{eq:v39-phase-equation}, use invariance, and write
$e^{iv\eta_R}=1+(e^{iv}-1)\eta_R$. Separating
$e^{i(u\cdot\kappa_R+b\tau_R)}-1$ gives
\eqref{eq:v40-residue-source-identity} with the displayed signs.
The final assertions follow directly. For representatives with
$|(u,b)|\le1$ its last term has modulus at most $C|(u,b)|$;
no division by $e^{iv}-1$ is uniform as $v\to0$.
\end{proof}

\begin{proposition}[The same Gaussian curvature at each resonance]
\label{prop:v40-resonant-curvature}
Put $\mu_R=(0,0,\bar\tau_R,c)$. On a coordinate neighborhood of a
resonance $\gamma$, the simple eigenvalue and projection are analytic
in the displacement $\delta\in\mathbb C^4$, and
\begin{equation}\label{eq:v40-resonant-expansion}
 \lambda_\gamma(\delta)
  =\zeta_\gamma\exp\left(i\delta\cdot\mu_R
          -\tfrac12\delta^{\mathsf T}\Omega_R\delta
          +O(|\delta|^3)\right),\qquad
 \Pi_\gamma(\delta)=\Pi_\gamma+O(|\delta|).
\end{equation}
The constants and radii may be chosen uniformly over all actual
resonances and $R\in I$. In particular the centered modulus has a
Gaussian upper bound for small real $\delta$.
\end{proposition}
\begin{proof}
Use the analytic family
$V_\delta=e^{-i\delta\cdot\mu_R}\mathscr T_{R,\gamma+\delta}$.
Its strong analyticity follows from the image-side roof and the
source-side occupation realization, not from a presumed bounded
forward-roof multiplier. Isolate the simple eigenvalue
$\zeta_\gamma$ by a resolvent contour. Let $l_\gamma$ be its left
functional with $l_\gamma(q_\gamma\nu)=1$. On bounded physical
strong vectors, \eqref{eq:v40-projection-formula} identifies it as
integration against $\overline q_\gamma$.

For a fixed real direction $w$ put $X=w\cdot X_R$ and differentiate
$V_{tw}$ in the scalar variable $t$ at zero. Denote its eigenvalue by
$\widetilde\lambda_\gamma(t)=e^{-itw\cdot\mu_R}\lambda_\gamma(tw)$. A first derivative, acting on a bounded physical
strong vector $h$, has physical representative
$i\mathscr T_{R,\gamma}(Xh)$. This is a description of the derivative
of the entire operator family; no boundedness of $M_X$ in isolation
is needed. In particular
\[
 V'_0q_\gamma=i\zeta_\gamma q_\gamma(X\circ T_R^{-1}),
 \qquad l_\gamma(V'_0q_\gamma)=0.
\]
The complementary powers on this centered vector are summable in the
strong norm. The derivative of the normalized right eigenvector is
therefore
\[
 e'_0=i\sum_{k\ge1}q_\gamma(X\circ T_R^{-k})
\]
with strong convergence; each partial sum is a bounded physical
representative. Applying $l_\gamma$ to the second differentiated
eigen-equation gives
\[
 \widetilde\lambda''_\gamma(0)
   =-\zeta_\gamma\left(\nu(X^2)
             +2\sum_{k\ge1}\nu(X(X\circ T_R^k))\right).
\]
Continuity of $l_\gamma V'_0$ justifies passing to the strong sum.
The physical covariance series is absolutely summable. Polarization
identifies the Hessian as $-\Omega_R$ for the logarithm of the
centered eigenvalue. Uniform contour analyticity bounds the cubic
remainder by Cauchy's estimate and gives the projector estimate.

For uniformity, the joint resonance set is compact by
Theorem~\ref{thm:v40-full-peripheral} and lies in $b=0$. All its
peripheral eigenvalues are simple and separated from complementary
spectrum locally. Strong-to-weak spectral stability and a finite
cover give common contour bounds on the local spaces. The uniform
positive lower bound for $\Omega_R$ then gives the Gaussian modulus
bound after reducing the coordinate radius. This argument does not
assert that the set of resonances of one radius persists at all
nearby radii.
\end{proof}
